\section{Problem 3: First Unique Employee}

\subsection{Problem Overview and Objectives}
A company records the access log sequence of $n$ employee names. An employee is defined as \textbf{unique} if their name occurs exactly once in the complete sequence.

The objectives are:
\begin{enumerate}
    \item Design \textbf{Approach 1}: Find the first unique employee name using \textbf{ONLY} the given list and repeated comparisons, using zero additional data structures ($O(1)$ auxiliary space).
    \item Design \textbf{Approach 2}: Find the first unique employee name using an additional data structure (\textbf{Hash Map / Frequency Map}) to optimize time complexity.
    \item Compare both solutions across time, space, and scalability for small vs. massive datasets.
\end{enumerate}

\subsection{Specifications}
\begin{itemize}
    \item \textbf{Input}: \texttt{employee\_sequence}, a list of $n$ string names ($K = \text{max string length}$).
    \item \textbf{Output}: The string name of the first unique employee in original sequence order, or \texttt{None} if all names have duplicates.
    \item \textbf{Constraints}: Approach 1 must allocate zero extra arrays, sets, or hash tables.
\end{itemize}

\subsection{Formal Algorithm Design}

\begin{algorithm}[H]
\caption{Approach 1: List-Only Repeated Comparisons ($O(1)$ Auxiliary Space)}
\begin{algorithmic}[1]
\Function{FindFirstUniqueListOnly}{$sequence$}
    \State $n \gets \text{Length}(sequence)$
    \For{$i \gets 0$ \textbf{to} $n - 1$}
        \State $target \gets sequence[i]$
        \State $is\_duplicate \gets \text{False}$
        \For{$j \gets 0$ \textbf{to} $n - 1$}
            \If{$i \neq j$ \textbf{and} \Call{CompareStrings}{$target, sequence[j]$} $== 0$}
                \State $is\_duplicate \gets \text{True}$
                \State \textbf{break} \Comment{Duplicate found, try next candidate}
            \EndIf
        \EndFor
        \If{\textbf{not} $is\_duplicate$}
            \State \Return $target$ \Comment{First unique employee found}
        \EndIf
    \EndFor
    \State \Return $\text{None}$
\EndFunction
\end{algorithmic}
\end{algorithm}

\begin{algorithm}[H]
\caption{Approach 2: Hash Map / Frequency Map ($O(n \cdot K)$ Time)}
\begin{algorithmic}[1]
\Function{FindFirstUniqueHashMap}{$sequence$}
    \State $freq\_map \gets \{\}$
    \For{\textbf{each} $name$ \textbf{in} $sequence$} \Comment{Pass 1: Count Frequencies}
        \State $freq\_map[name] \gets freq\_map.\text{get}(name, 0) + 1$
    \EndFor
    \For{\textbf{each} $name$ \textbf{in} $sequence$} \Comment{Pass 2: Preserve Original Sequence Order}
        \If{$freq\_map[name] == 1$}
            \State \Return $name$ \Comment{First unique employee found}
        \EndIf
    \EndFor
    \State \Return $\text{None}$
\EndFunction
\end{algorithmic}
\end{algorithm}

\subsection{Complexity Analysis and Comparison}

\subsubsection{Approach 1 (List-Only Repeated Comparisons)}
\begin{itemize}
    \item \textbf{Best-Case Time Complexity}: $O(n \cdot K)$.
    Occurs when the first element $sequence[0]$ is unique. It compares $sequence[0]$ against all $n-1$ other elements, taking $(n-1) \times O(K) = O(n \cdot K)$ operations.
    
    \item \textbf{Worst-Case Time Complexity}: $O(n^2 \cdot K)$.
    Occurs when no unique element exists, or the unique element is located at the end. Outer loop runs $n$ times, inner loop runs up to $n$ times, yielding $n \times n \times O(K) = O(n^2 \cdot K)$.
    
    \item \textbf{Auxiliary Space Complexity}: $O(1)$.
    Uses zero additional memory data structures; operates entirely within the input list.
\end{itemize}

\subsubsection{Approach 2 (Hash Map / Frequency Map)}
\begin{itemize}
    \item \textbf{Best, Average, \& Worst-Case Time Complexity}: $O(n \cdot K)$.
    Pass 1 inserts $n$ names into the hash map, performing string hashing and key matching in $n \times O(K)$ time. Pass 2 iterates $n$ names with $O(K)$ hash lookups. Total Time $= 2 \cdot n \cdot O(K) = O(n \cdot K)$.
    
    \item \textbf{Auxiliary Space Complexity}: $O(n \cdot K)$.
    Stores up to $n$ distinct name key strings of max length $K$ in the hash map.
\end{itemize}

\begin{table}[htbp]
\centering
\small
\caption{Comparison Matrix for First Unique Employee Strategies}
\label{tab:problem3_complexity}
\begin{tabularx}{\textwidth}{@{} >{\raggedright\arraybackslash}X c c c c @{}}
\toprule
\textbf{Approach} & \textbf{Best Time} & \textbf{Average Time} & \textbf{Worst Time} & \textbf{Auxiliary Space} \\
\midrule
\textbf{1. List-Only Comparisons} & $O(n \cdot K)$ & $O(n^2 \cdot K)$ & $O(n^2 \cdot K)$ & $O(1)$ \\
\textbf{2. Hash Map (Frequency Map)} & $O(n \cdot K)$ & $O(n \cdot K)$ & $O(n \cdot K)$ & $O(n \cdot K)$ \\
\bottomrule
\end{tabularx}
\end{table}

\subsection{Scalability and Selection Justification}
\begin{itemize}
    \item \textbf{Small Datasets} ($n \le 100$): Approach 1 is lightweight and requires zero memory allocations.
    \item \textbf{Large Datasets} ($n \ge 100,000$): Approach 1 becomes impractical due to quadratic $O(n^2 \cdot K)$ growth ($\approx 10^{10}$ comparisons). Approach 2 completes in milliseconds ($\approx 2 \cdot 10^5$ operations). Approach 2 is strongly preferred for production system logs.
\end{itemize}
